\chapter{INTRODUCTION}

\section{Background}
In this modern age,  rapid increase of population, climate change, degradation in the quality of soil, and declining availability of arable land have placed significant pressure on conventional agricultural practices. Considering this fact, hydroponic farming is used to cultivate crops in vertically stacked layers instead of traditional horizontal farming methods.\\[10pt]
In hydroponics, we externally provide nutrients that plants require to grow and develop. Nutrients needed for the plant are mixed within the water and thus directly provide the nutrients to the plant roots. All plants will get an equal share of nutrients, which leads all plants to grow equally.
Hydroponics started to flourish in the early twentieth century, when William Frederick Gericke of the University of California popularized the idea of growing crops in nutrient solutions rather than soil for large-scale production. The reason behind hydroponic farming is to maximize production in limited spaces, higher crop yield, and faster growth.\\[10pt]
As technology has become vital and more accessible, our project aims to automate hydroponic systems where temperature, humidity, pH, and EC can be monitored and controlled. These parameters are essential for maintaining optimal growing conditions and ensuring healthy plant growth. Among all the parameters, pH and EC are prioritized as plant growth depends directly on the amount of micro-nutrients and macro-nutrients which EC and pH directly controls the availability and absorption of the nutrients. The integration of Internet of Things (IoT) can minimize human operations as automated nutrient delivery using sensors, controllers, and actuators to continuously monitor and regulate nutrient levels in real time. The collected data is processed by a microcontroller, which determines whether the nutrient concentration is within the desired range for healthy plant growth. Growers can track nutrient levels, compensation efficiency, and plant conditions through web-based interfaces.\\[10pt]
In this context, the project focuses on the design and implementation of an automated hydroponic system in order to provide accurate nutrient delivery. The designed approach minimizes human error, reduces nutrient wastage, and sustains small-scale indoor agricultural applications.

\section{Motivation}
In a densely populated city facing rapid urbanization, such as Kathmandu, the toll taken on traditional farming has been massive. With less and less arable land, water scarcity, and lower crop yield with time, conventional soil-based farming is facing increasing challenges and turning into a struggle to keep up with. Hydroponics mitigates such drawbacks since it does not require a big space, does not require soil, is free from the majority of disease-causing parasites, circulates required amounts of water and nutrients, thus promoting water recycling and higher yield. Seeing this, we decided that the best method that upgrades current farming methods with minimal drawbacks while also being highly automatable was an automatic nutrient-dosing hydroponic system.\\[10pt]
Hydroponics, even done manually, is a highly effective and sustainable approach to cultivation. Introducing automation to it reduces time losses, cuts down human errors, and boosts the accuracy, precision, and overall performance of the model. This ultimately leads to high crop yield, high production, and less water use, all with minimal effort and fewer challenges through smart management of resources.

\section{Problem Definition}
Hydroponics requires constant monitoring and utmost precision, which is difficult to achieve through manual labour, but is generally done manually. Parameters such as pH, EC, and humidity differ from plant to plant, which must be provided with a minimum tolerance. Even small deviations, if unnoticed for a longer time, which is very likely with manual models, can cause nutrient deficiencies, root stress, and ultimately crop loss. Furthermore, the manual approach lacks data logs and real-time feedback, which is essential for improving overall efficiency. All these combined instances pose a threat to crop production in hydroponics, which requires optimal performance and efficient resource utilization. Even though automatic approaches have been introduced and implemented for some time, existing hydroponic automation often lacks cost-effective control strategies, real-time feedback mechanisms, and stability-aware dosing approaches suitable for low-cost embedded platforms.

\section{Objectives}
\begin{itemize}
    \item To design and implement a fully automated closed-loop hydroponic system.
    \item To ensure precise real-time monitoring and remote control of key parameters.
\end{itemize}

\section{Scope and Application}
\begin{itemize}
    \item During experimental validation, the system is configured for a single crop i.e. lettuce at a time to ensure precise control and reliable performance evaluation.
    \item Provides accurate nutrient management for each individual crop through a fully automated hydroponic system with closed-loop pH and EC control.
    \item Sensors, automated dosing pumps, and microcontroller-based control are integrated to minimize human interaction, lowering error rates, and increasing system reliability.
    \item Applicable in settings with scarcity of land and water resources.
    \item It is applicable for home-based systems, educational institutions and small commercial farms due to its affordable and scalable architecture.
    \item Large-scale commercial hydroponic installations with multi-crop and industrial-grade control requirements are considered outside the scope of this project.
    \item Achievement of sustainable farming through increased crop yield and effective use of water and nutrients.
\end{itemize}

\section{Report Organization}
This report consists of eight main chapters:

\subsection{Introduction}
The Introduction chapter provides an overview of the project, which includes background systems, motivation behind the project, problem definition, objectives, and scope of the proposed system. 

\subsection{Literature Review}
This chapter reviews existing research and developments in hydroponic systems, automation techniques, and IoT-based monitoring solutions, highlighting their advantages and limitations.

\subsection{Requirement Analysis}
The chapter provides the functional and non-functional requirements of the system, as well as hardware requirements and constraints considered during design.

\subsection{Methodology}
This chapter explains the overall system design, including system architecture, sensor subsystem, control strategies, actuation mechanisms, user interface, and working principles of the system.

\subsection{Implementation Details}
This chapter describes the hardware and software implementation of the system, including sensor interfacing, ESP32 firmware, backend, database, and frontend development.

\subsection{Result and Analysis}
This chapter presents the experimental results obtained from the system such as sensor data visualization, system performance evaluation, and observed issues.

\subsection{Future Enhancements}
This chapter discusses possible improvements such as multi-crop support, machine learning integration, enhanced calibration techniques, and mobile application development.

\subsection{Conclusion }
This chapter summarizes the overall work, key achievements, and final outcomes of the project. 
Finally, the report includes Appendices (budget, schedule, circuit diagrams, PCB design) and References used throughout the project.
